\section{Análisis cualitativo de fortalezas y debilidades}
\label{sec:cualitativo}

\subsection{Clasificación de las respuestas}
Para identificar el origen de los fallos se clasificaron las 7.680 respuestas de
la evaluación en la partición de prueba (sección~\ref{sec:prueba}) en categorías
excluyentes, asignadas en el orden de la tabla~\ref{tab:errores}: primero las
correctas (\emph{reward} 1); luego las que alcanzaron el límite de tokens; las
que no cumplen el contrato de formato; las que alteran alguna pista original, y,
por último, las grillas con el formato correcto y las pistas intactas que no
resuelven el puzzle. La clasificación usa el motivo de término y los componentes
del \emph{reward} registrados en cada respuesta.
Además, se volvió a evaluar cada una de las 7.680 respuestas con el verificador
actual: todos los \emph{rewards} recalculados coinciden con los registrados.

\begin{table}[htbp]
  \centering\small
  \caption{Clasificación de las respuestas de la prueba por modelo, \emph{shots} y
  tamaño de tablero (960 respuestas por fila: 60 tareas por 16 respuestas).}
  \label{tab:errores}
  \begin{tabularx}{\linewidth}{@{}lccYYYYY@{}}
    \toprule
    Modelo & \emph{Shots} & Tablero & Correcta & Truncada & Formato
      & Altera pistas & Grilla incorrecta \\
    \midrule
    2B & 0 & 4x4 & 529 & 345 & 81 & 5 & 0 \\
       & 0 & 6x6 &  67 & 889 &  4 & 0 & 0 \\
       & 2 & 4x4 & 621 & 301 &  0 & 28 & 10 \\
       & 2 & 6x6 & 137 & 822 &  0 & 1 & 0 \\
    \midrule
    9B & 0 & 4x4 & 856 & 104 &  0 & 0 & 0 \\
       & 0 & 6x6 & 287 & 672 &  1 & 0 & 0 \\
       & 2 & 4x4 & 855 &  97 &  0 & 3 & 5 \\
       & 2 & 6x6 & 329 & 615 &  0 & 0 & 16 \\
    \bottomrule
  \end{tabularx}
\end{table}

El fallo dominante es el truncamiento: 3.851 de las 7.680 respuestas (50,1\,\%)
alcanzaron el límite de 2.304 tokens, el 61,4\,\% de las del 2B y el 38,8\,\% de
las del 9B, y en 6x6 es la causa de entre el 64\,\% y el 93\,\% de las respuestas
según el modelo y la configuración. Casi todos los demás fallos son de formato o
de pistas alteradas; las grillas con el formato correcto y las pistas intactas
pero incorrectas son raras (31 respuestas, el 0,4\,\%). En ellas, con 2 \emph{shots}
y en 4x4, el 2B rellena en la mediana el 44\,\% de las casillas vacías con el valor
correcto (\emph{fill} del \emph{reward}). Por lo tanto, en este régimen de
razonamiento largo el modelo rara vez entrega una grilla completa y errónea:
resuelve el puzzle o no termina de escribir su razonamiento, y es el presupuesto
de tokens el que decide cuál de los dos ocurre. Esta clasificación no distingue
cuántas respuestas truncadas habrían llegado a una solución correcta con más
tokens; en el piloto 2 (sección~\ref{sec:pilotos}), el 9B resolvió las tres tareas de 6x6
probadas (una por nivel) con 8.192 tokens, aunque con tan pocas tareas es solo
una indicación.

\subsection{Ejemplo de acierto y fortaleza observada}
Los ejemplos de esta sección se eligieron para ilustrar las categorías
de la tabla~\ref{tab:errores}; no son una muestra aleatoria.

\textbf{Modelo, nivel e instancia:} Qwen3.5-2B, 4x4 con 4 casillas vacías, 2
\emph{shots}, \texttt{test-4x4-n4-00000}, respuesta 0 de 16.

\textbf{Problema y respuesta esperada:}
\[
  \text{Puzzle: }
  \begin{bmatrix} 0&3&2&0\\ 2&4&1&3\\ 4&2&3&1\\ 0&1&4&0 \end{bmatrix}
  \qquad
  \text{Solución: }
  \begin{bmatrix} 1&3&2&4\\ 2&4&1&3\\ 4&2&3&1\\ 3&1&4&2 \end{bmatrix}
\]

\textbf{Respuesta del modelo y resultado del verificador:} el modelo razona por
filas (\emph{Step 1: Analyze Row 1\ldots}), deduce las casillas faltantes a partir
de las columnas, comprueba los cuatro bloques y termina con una grilla idéntica a
la solución, en el formato requerido y con 605 tokens. \emph{Reward} 1.

Con pocas casillas vacías y ejemplos en el prompt, el 2B aplica correctamente la
deducción por casilla única: en esta tarea acierta 15 de las 16 respuestas. El
acierto, sin embargo, no es estable en todas las tareas: en la siguiente instancia
(\texttt{test-4x4-n4-00001}) el 2B falla de la forma que se describe más abajo, y
la respuesta que falla en esta misma tarea es la del segundo ejemplo de error. Esta
variabilidad entre intentos es la que miden \passuno\ y \passk, y es la condición
que permite que el aprendizaje reforzado tenga señal.

\subsection{Ejemplos de error y debilidades observadas}

\paragraph{Respuesta correcta rechazada por formato.}
\textbf{Modelo, nivel e instancia:} Qwen3.5-2B, 4x4 con 4 casillas vacías, 0
\emph{shots}, \texttt{test-4x4-n4-00000}, respuesta 12 de 16.

\textbf{Respuesta del modelo y resultado del verificador:} tras 889 tokens de
análisis, el modelo escribe la grilla correcta, pero dentro del bloque
\texttt{<answer>} como filas de texto (\texttt{1 3 2 4}, \texttt{2 4 1 3}, \ldots)
y no como matriz JSON. El contrato exige la matriz JSON, por lo que el verificador
la rechaza: \emph{reward} 0.

\textbf{Origen del error:} formato. El modelo resolvió la tarea, pero no respetó
el contrato de salida. Es el fallo más frecuente del 2B sin ejemplos en 4x4 (81
de 960 respuestas; en 69 de ellas el bloque contiene filas de texto y en 12 no
hay bloque de respuesta) y no aparece con 2 \emph{shots} (0 de 960), cuyos
ejemplos muestran la matriz JSON. Es el costo de exigir un contrato estricto
(sección~\ref{sec:robustez}).

\paragraph{Pista original alterada.}
\textbf{Modelo, nivel e instancia:} Qwen3.5-2B, 4x4 con 4 casillas vacías, 2
\emph{shots}, \texttt{test-4x4-n4-00000}, respuesta 8 de 16.

\textbf{Respuesta del modelo y resultado del verificador:} el modelo completa
las casillas paso a paso, pero en la última fila escribe $[4,1,2,3]$ en lugar
de $[3,1,4,2]$: cambia la pista 4 de la columna 3 por un 2 y coloca valores
erróneos en las casillas vacías. Solo la mitad de las casillas vacías coincide con la solución y 11 de
las 12 pistas se conservan: \emph{reward} 0.

\textbf{Origen del error:} razonamiento y copia. En 4x4 con 2 \emph{shots}, 28
respuestas del 2B alteran alguna pista; con 0 \emph{shots}, 5.

\paragraph{Puzzle declarado ya resuelto.}
\textbf{Modelo, nivel e instancia:} Qwen3.5-2B, 4x4 con 4 casillas vacías, 2
\emph{shots}, \texttt{test-4x4-n4-00001}, respuesta 5 de 16.

\textbf{Respuesta del modelo y resultado del verificador:} el modelo verifica las
filas y columnas, afirma que ``el puzzle ya está completamente resuelto'' y
devuelve la grilla sin cambios, con 4 casillas en 0. El formato es correcto y las
pistas se conservan, pero ninguna casilla vacía se completa (\emph{fill} 0):
\emph{reward} 0.

\textbf{Origen del error:} razonamiento. El modelo no detecta que las casillas
con 0 están vacías. Es uno de los pocos casos de grilla con el formato correcto
y pistas intactas, la categoría menos frecuente de la tabla~\ref{tab:errores}.

\paragraph{Razonamiento truncado.}
\textbf{Modelo, nivel e instancia:} Qwen3.5-2B, 6x6 con 16 casillas vacías, 0
\emph{shots}, \texttt{test-6x6-n16-00000}, respuesta 0 de 16.

\textbf{Respuesta del modelo y resultado del verificador:} el modelo analiza fila
por fila, intersecta conjuntos de números faltantes y vuelve a las mismas filas,
hasta agotar los 2.304 tokens en medio de una frase, sin llegar a escribir el bloque
\texttt{<answer>}. \emph{Reward} 0.

\textbf{Origen del error:} presupuesto de tokens y razonamiento a la vez. Es el
comportamiento de la mayoría de las respuestas de 6x6 (tabla~\ref{tab:errores}):
el modelo deduce casilla por casilla, pero el razonamiento en lenguaje natural no
cabe en el límite.
